随着高校毕业论文及各类学术论文写作任务的增多，论文排版的规范性与编辑效率逐渐受到关注。传统文字处理工具虽然操作直观，但在复杂公式、图表编号、交叉引用、参考文献管理和模板适配等方面仍需大量人工调整；LaTeX 虽然具备较强的学术排版能力，但其命令语法、环境结构、编译配置和错误日志分析对初学者并不友好。针对普通用户使用 LaTeX 时存在的代码记忆成本高、模板使用复杂、编译错误难以定位等问题，本文设计并实现了一套基于 JavaWeb 和 LaTeX 的学术论文智能排版系统，为论文写作与排版提供在线化、组件化和智能化支持。 

系统采用前后端分离架构，前端基于 Vue 3 和 Vite 实现页面展示与交互功能，后端基于 Spring Boot、MyBatis 和 MySQL 完成业务逻辑处理与数据持久化。系统实现了用户注册登录、论文项目管理、LaTeX 在线编辑、模板管理、论文编译、PDF 预览、文件导出和 AI 辅助等功能。用户可在浏览器中创建论文项目，选择并导入论文模板，在线编辑论文内容，并通过不同编译引擎生成 PDF 文档。 

系统的核心设计是面向论文结构的组件库拖拽式插入机制。系统将章节、图片、表格、公式、符号等常用 LaTeX 结构封装为可视化组件，用户无需专门记忆完整的 LaTeX 命令和环境写法，只需拖入所需组件，系统即可自动生成对应代码框架。该方式降低了 LaTeX 语法使用门槛，减少了命令拼写错误、括号缺失和环境未闭合等问题，提高了论文编辑效率。 